\section{Physical contractions and Gaussian source estimates}

We retain the original source positions, local branch labels, and common
Gaussian law of Definition~\ref{def:peps_source_gaussian_law}.
The estimates below concern actual allowed words and specified isometric
input frames. They allow different ket and bra input spaces.
Finite orthonormal coordinates are chosen whenever matrices or partial
traces are used; their dimensions are unrestricted. This follows the
finite-space convention in \cite[proof of Theorem~5.2]{OpenAI2026PolynomialPEPS},
\texttt{04-compression.tex}, lines 5--19 and 383--558.
The intrinsic component and Gram maps also make sense for arbitrary
inner-product discard spaces.

\begin{definition}[Physical components of an actual word]
    \label{def:peps_physical_component}
    \lean{TNLean.PEPS.PairEffect.Word.physicalComponent}
    \lean{TNLean.PEPS.PairEffect.Word.physicalComponentCoordinates}
    \lean{TNLean.PEPS.PairEffect.Layout.physicalDiscardIso}
    \lean{TNLean.PEPS.PairEffect.Layout.physicalDiscardIso_tmul}
    \leanok
    Let $V:H\to K$ be the evaluation of a word and let
    $J:K\simeq\C^X\otimes D$ be a linear isometric equivalence. For $x\in X$, put
    \begin{align}
        C_{x,V}=(\langle x|\otimes I_D)JV:H\longrightarrow D.
    \end{align}
    After choosing finite orthonormal bases of $H$ and $D$, let $A_{x,V}$
    be its matrix. For an output layout consisting of physical registers
    followed by discarded registers, the canonical choice of $J$ is the
    concatenation isometry followed by the physical coordinate isometry.
    It sends a concatenated elementary tensor $s\otimes t$ to
    $[s]_{\rm phys}\otimes t$.
\end{definition}

\begin{theorem}[Component contraction]
    \label{thm:peps_physical_component_contraction}
    \lean{TNLean.PEPS.PairEffect.Word.norm_physicalComponent_le_one}
    \lean{TNLean.PEPS.PairEffect.Word.norm_physicalComponentCoordinates_le_one}
    \leanok
    \uses{def:peps_physical_component}
    If $V$ is allowed, then $\|C_{x,V}\|\leq1$ and $\|A_{x,V}\|\leq1$.
    The first bound requires neither completeness nor finite dimension of $D$.
\end{theorem}

\begin{proof}\leanok
    \uses{def:peps_physical_component}
    Since $C_{x,V}=(\langle x|\otimes I_D)JV$, the coordinate
    functional, the isometry $J$ and the allowed word give
    $\|C_{x,V}\|\leq\|\langle x|\otimes I_D\|\,\|J\|\,\|V\|\leq1$.
    Orthonormal coordinates preserve this bound.
\end{proof}

\begin{definition}[Mixed physical Gram operator]
    \label{def:peps_physical_gram}
    \lean{TNLean.PEPS.PairEffect.Word.physicalGramMap}
    \lean{TNLean.PEPS.PairEffect.Word.physicalGramMap_inner}
    \lean{TNLean.PEPS.PairEffect.Word.physicalGramMatrix}
    \lean{TNLean.PEPS.PairEffect.Word.physicalGramMatrix_apply}
    \leanok
    Let $W:H'\to K$ be another word with the same output identification $J$.
    Assume that $H,H'$ have finite orthonormal bases. Define $O:H\to H'$ by
    \begin{align}
        \langle t,Os\rangle=\langle C_{y,W}t,C_{x,V}s\rangle.
    \end{align}
    Its matrix has bra rows and ket columns:
    \begin{align}
        O_{ji}=\langle C_{y,W}e'_j,C_{x,V}e_i\rangle.
    \end{align}
    Only the finite-dimensional domain of the Riesz representation is needed;
    the discard space $D$ may be incomplete.
\end{definition}

\begin{theorem}[Gram contraction]
    \label{thm:peps_physical_gram_contraction}
    \lean{TNLean.PEPS.PairEffect.Word.norm_physicalGramMap_le_one}
    \lean{TNLean.PEPS.PairEffect.Word.norm_physicalGramMatrix_le_one}
    \leanok
    \uses{def:peps_physical_gram,thm:peps_physical_component_contraction}
    For allowed $V,W$, both the mixed Gram operator and its matrix satisfy
    $\|O\|\leq1$.
\end{theorem}

\begin{proof}\leanok
    \uses{thm:peps_physical_component_contraction}
    By Cauchy--Schwarz and
    Theorem~\ref{thm:peps_physical_component_contraction},
    $|\langle t,Os\rangle|
        =|\langle C_{y,W}t,C_{x,V}s\rangle|
        \leq\|C_{y,W}t\|\,\|C_{x,V}s\|
        \leq\|t\|\,\|s\|$.
    Passing to orthonormal coordinates preserves this estimate.
\end{proof}

\begin{definition}[Joint physical and discarded output coordinates]
    \label{def:peps_physical_output_matrix}
    \lean{TNLean.PEPS.PairEffect.Word.physicalOutputMatrix}
    \leanok
    \uses{def:peps_physical_component}
    Choose a finite orthonormal basis of $D$. Write $K_V$ for the matrix of $JV$
    in the resulting product output basis and the chosen input basis.
    Its rows are indexed by a physical coordinate and a discard coordinate.
\end{definition}

\begin{theorem}[Actual partial trace and exterior contraction]
    \label{thm:peps_exterior_contraction}
    \lean{TNLean.PEPS.PairEffect.Word.partialTrace_physicalOutputMatrix_mul_conjTranspose}
    \lean{TNLean.PEPS.PairEffect.Word.norm_physicalOutputMatrix_le_one}
    \lean{TNLean.PEPS.PairEffect.Word.rectangularTraceNorm_partialTrace_physicalOutputMatrix_le}
    \leanok
    \uses{def:peps_physical_output_matrix,def:peps_physical_gram}
    For an arbitrary rectangular matrix $\rho:H'\to H$,
    \begin{align}
        \bigl[\operatorname{tr}_{D}(K_V\rho K_W^\dagger)\bigr]_{xy}
         & =\sum_{i,j}\rho_{ij}O_{ji}.
    \end{align}
    If $V,W$ are allowed, then $\|K_V\|,\|K_W\|\leq1$ and
    \begin{align}
        \|\operatorname{tr}_{D}(K_V\rho K_W^\dagger)\|_1\leq\|\rho\|_1.
    \end{align}
    No positivity or self-adjointness is required of $\rho$.
\end{theorem}

\begin{proof}\leanok
    Expand the partial trace in the discard basis and use the Gram
    entry formula. The output matrices are contractions in orthonormal
    coordinates. The rectangular trace-norm contraction estimate for their
    mixed partial trace gives the last assertion.
\end{proof}

\begin{definition}[Proper input frames]
    \label{def:peps_physical_frames}
    \lean{TNLean.PEPS.PairEffect.Word.physicalComponentFrameMatrix}
    \lean{TNLean.PEPS.PairEffect.Word.physicalComponentFrameMatrix_apply}
    \lean{TNLean.PEPS.PairEffect.Word.physicalFrameGramMatrix}
    \lean{TNLean.PEPS.PairEffect.Word.physicalFrameGramMatrix_apply}
    \lean{TNLean.PEPS.PairEffect.Word.physicalOutputFrameMatrix}
    \leanok
    \uses{def:peps_physical_component,def:peps_physical_gram}
    Let $F:\C^n\to H$ and $F':\C^{n'}\to H'$ be linear isometries.
    Replace the orthonormal input coordinate maps of
    Definitions~\ref{def:peps_physical_component} and~\ref{def:peps_physical_gram}
    by $F,F'$, respectively. Thus
    \begin{align}
        A_{x,V,F} & =[C_{x,V}F],                               &
        O_{ji}    & =\langle C_{y,W}F'e'_j,C_{x,V}Fe_i\rangle, &
        K_{V,F}   & =[JVF].
    \end{align}
    These frames need not span either input space.
\end{definition}

\begin{theorem}[Contraction in proper frames]
    \label{thm:peps_physical_frame_contraction}
    \lean{TNLean.PEPS.PairEffect.Word.norm_physicalComponentFrameMatrix_le_one}
    \lean{TNLean.PEPS.PairEffect.Word.norm_physicalFrameGramMatrix_le_one}
    \lean{TNLean.PEPS.PairEffect.Word.norm_physicalOutputFrameMatrix_le_one}
    \lean{TNLean.PEPS.PairEffect.Word.rectangularTraceNorm_partialTrace_physicalOutputFrameMatrix_le}
    \leanok
    \uses{def:peps_physical_frames,thm:peps_physical_component_contraction}
    For allowed words, the three matrices in
    Definition~\ref{def:peps_physical_frames} have operator norm at most one.
    For every rectangular $Z$,
    \begin{align}
        \|\operatorname{tr}_{D}(K_{V,F}ZK_{W,F'}^\dagger)\|_1\leq\|Z\|_1.
    \end{align}
\end{theorem}

\begin{proof}\leanok
    Compose the actual word with the given isometric input frame.
    The component and mixed Gram estimates, followed by the rectangular
    partial-trace inequality, apply without a spanning condition.
\end{proof}

\begin{definition}[Schmidt-weighted physical vectors]
    \label{def:peps_schmidt_physical_vector}
    \lean{TNLean.PEPS.PairEffect.Word.schmidtPhysicalVector}
    \lean{TNLean.PEPS.PairEffect.Word.schmidtPhysicalVector_repr}
    \leanok
    \uses{def:peps_physical_frames}
    Let $F:\C^{I\times T}\to H$ be an isometry and let $\tau_t\in\R$.
    For fixed physical coordinate $x$, set
    \begin{align}
        \psi_i=\sum_{t\in T}\sqrt{\tau_t}\,|t\rangle\otimes C_{x,V}F|i,t\rangle.
    \end{align}
    In the product basis its $(t,d)$ coordinate is
    $\sqrt{\tau_t}(A_{x,V,F})_{d,(i,t)}$.
    A bra word, independent index sets $L,U$, frame $F'$, and weights $\tau'$
    define vectors $\phi_l$ in the same manner.
\end{definition}

\begin{theorem}[Weighted source matrix from actual vectors]
    \label{thm:peps_weighted_physical_source_identity}
    \lean{TNLean.PEPS.PairEffect.Word.sum_schmidtPhysicalVector_outer_eq_weightedSourceError}
    \leanok
    \uses{def:peps_schmidt_physical_vector}
    For arbitrary complex coefficients $z_{il}$, tracing the common discard
    factor in $\sum_{i,l}z_{il}|\psi_i\rangle\langle\phi_l|$ gives the
    rectangular matrix $W:T\leftarrow U$ with entries
    \begin{align}
        W_{tu}=\sqrt{\tau_t\tau'_u}\sum_{i,l}z_{il}O_{(l,u),(i,t)}
    \end{align}
    when both weight families are nonnegative. For arbitrary real weights,
    the factor is $\sqrt{\tau_t}\sqrt{\tau'_u}$.
    The coefficient $z_{il}$ is not conjugated. The bra coordinates in the
    outer product supply the conjugation in $O$.
\end{theorem}

\begin{proof}\leanok
    Expand both actual vectors in their product bases and sum over the
    common discard coordinate. This gives precisely the weighted source matrix
    with the ordinary transpose of the bra-row, ket-column Gram matrix.
\end{proof}

\begin{definition}[Separated output and its physical operator]
    \label{def:peps_separated_physical_density}
    \lean{TNLean.PEPS.PairEffect.Word.separatedSchmidtOutput}
    \lean{TNLean.PEPS.PairEffect.Word.separatedSchmidtOutput_repr}
    \lean{TNLean.PEPS.PairEffect.Word.separatedPhysicalCoordinates}
    \lean{TNLean.PEPS.PairEffect.Word.separatedPhysicalDensity}
    \leanok
    \uses{def:peps_schmidt_physical_vector}
    Let $V':H_E\to K_E$ be an exterior word, let
    $J':K_E\simeq\C^Y\otimes D'$ be its output identification, and let
    $G:\C^T\to H_E$ be an isometry.
    The joint output at coordinate $i$ is
    \begin{align}
        \Psi_i=\sum_t\sqrt{\tau_t}\,(JVF|i,t\rangle)\otimes(J'V'G|t\rangle).
    \end{align}
    A separate bra word and its own frame and weights define $\Phi_l$.
    Reorder the four coordinate indices so that physical coordinates precede
    discarded coordinates, and put
    \begin{align}
        M(z)=\operatorname{tr}_{D\otimes D'}
        \sum_{i,l}z_{il}|\Psi_i\rangle\langle\Phi_l|.
    \end{align}
    At fixed affected physical coordinates $x,y$, the vectors used in this
    formula are those of Definition~\ref{def:peps_schmidt_physical_vector}
    tensored with the corresponding exterior outputs.
\end{definition}

\begin{theorem}[Tracing both discarded factors]
    \label{thm:peps_separated_physical_density_blocks}
    \lean{TNLean.PEPS.PairEffect.Word.partialTrace_separatedSchmidtOutput_eq_weightedSourceError}
    \lean{TNLean.PEPS.PairEffect.Word.separatedPhysicalDensity_submatrix}
    \leanok
    \uses{def:peps_separated_physical_density,thm:peps_weighted_physical_source_identity}
    For each pair $x,y$ of affected physical coordinates, the corresponding
    $Y\times Y$ block of $M(z)$ equals
    \begin{align}
        M(z)_{xy}=\operatorname{tr}_{D'}(K_{V',G}W_{xy}(z)K_{W',G'}^\dagger).
        \label{eq:peps_separated_physical_density_block}
    \end{align}
    This equality requires no allowedness assumption on the words and no
    normalization of the real weight families.
\end{theorem}

\begin{proof}\leanok
    Expand the finite Schmidt sums and both discard traces. The trace
    over $D$ is the weighted source matrix from
    Theorem~\ref{thm:peps_weighted_physical_source_identity}; the remaining
    trace gives~(\ref{eq:peps_separated_physical_density_block}).
\end{proof}

\begin{theorem}[Expected exterior trace norm under diagonal covariance]
    \label{thm:peps_weighted_physical_source_bound}
    \lean{TNLean.PEPS.PairEffect.Word.integral_rectangularTraceNorm_weightedPhysicalSource_le}
    \leanok
    \uses{thm:peps_physical_frame_contraction,def:peps_separated_physical_density}
    Assume all four words are allowed. Let $p_i,p'_l,\tau_t,\tau'_u$ be
    nonnegative probability weights and let $\kappa\geq0$. On a probability
    space, suppose every $z_{il}\overline{z_{i'l'}}$ is integrable and
    \begin{align}
        \mathbb E[z_{il}\overline{z_{i'l'}}]
        =\mathbf1_{i=i'}\mathbf1_{l=l'}\,
        \kappa\sqrt{p_i p'_l}.
    \end{align}
    Then, for each $x,y$, the exterior block satisfies
    \begin{align}
        \mathbb E\|M(z)_{xy}\|_1\leq\sqrt\kappa.
    \end{align}
    The input frames are only required to be isometries. The Gram contraction
    in this estimate follows from the actual allowed words.
\end{theorem}

\begin{proof}\leanok
    Use the Gram norm bound and the weighted-source second-moment
    estimate. The exterior partial trace is a trace-norm contraction.
    The pair-product integrability supplies the integrability needed to
    integrate this pointwise inequality.
\end{proof}

\begin{definition}[Original corrected Schmidt coordinates]
    \label{def:peps_original_schmidt_coordinates}
    \lean{TNLean.PEPS.PairEffect.SourceCircuit.CorrectedSchmidtCoordinates}
    \leanok
    \uses{def:peps_source_coordinates}
    For $S\subseteq\mathcal E$, write
    $\mathcal I_S=\prod_{e\in S}\{0,\ldots,r_e-1\}^2$.
    Each original position contributes two independent endpoint indices.
    The two local labels at each position are fixed before sampling.
\end{definition}

\begin{theorem}[Gaussian estimates for separated actual words]
    \label{thm:peps_gaussian_separated_physical_bound}
    \lean{TNLean.PEPS.PairEffect.SourceCircuit.integral_rectangularTraceNorm_gaussianPhysicalSource_le}
    \lean{TNLean.PEPS.PairEffect.SourceCircuit.integrable_gaussianSeparatedPhysicalDensity_and_integral_le}
    \leanok
    \uses{def:peps_original_schmidt_coordinates,thm:peps_actual_gaussian_covariance,thm:peps_weighted_physical_source_bound}
    Use the actual Gaussian coefficient product $Z_{u,v}$ at original positions
    $S$ and fixed local ket and bra labels. Assume $k>0$ and that the local
    Schmidt weights and the two crossing weight families are probability
    weights. Let the specified input frames be isometries and all four
    separated words be allowed. Then
    \begin{align}
        \mathbb E_k\|M(Z)_{xy}\|_1 & \leq k^{-|S|/2},       \\
        \mathbb E_k\|M(Z)\|_1      & \leq |X|^2 k^{-|S|/2}.
    \end{align}
    The full trace norm is integrable.
    Neither estimate depends on the exterior physical dimension or either
    discard dimension.
\end{theorem}

\begin{proof}\leanok
    The original-source law gives diagonal covariance with
    $\kappa=k^{-|S|}$ and product probabilities of total mass one.
    Apply Theorem~\ref{thm:peps_weighted_physical_source_bound} to each block.
    The trace norm of the full matrix is at most the sum of the block trace
    norms; sum over the $|X|^2$ affected coordinate pairs. Integrability follows
    from the finite Gaussian coefficient expansion.
\end{proof}

\begin{theorem}[Prepared matrices on mixed pure inputs]
    \label{thm:peps_prepared_vector_density}
    \lean{TNLean.PEPS.PairEffect.Word.preparedMatrix_mulVec_repr}
    \lean{TNLean.PEPS.PairEffect.Word.preparedMatrix_mul_vecMulVec}
    \lean{TNLean.PEPS.PairEffect.Word.preparedMatrix_mul_vecMulVec_preparedInput}
    \lean{TNLean.PEPS.PairEffect.SourceCircuit.sourceBasisMatrix_mul_vecMulVec}
    \leanok
    For the actual preparation followed by a word, let $T$ denote its linear
    map and $K_T$ its matrix in finite orthonormal bases. Then
    \begin{align}
        K_T[\psi] & =[T\psi],                      &
        K_T|\psi\rangle\langle\phi|K_{T'}^\dagger
                  & =|T\psi\rangle\langle T'\phi|.
    \end{align}
    These identities hold for arbitrary prepared source vectors, including
    the canonical source-basis substitutions in the original circuit. They
    also hold with preparation written explicitly on the input side of the
    residual word. No normalization or allowedness is required.
\end{theorem}

\begin{proof}\leanok
    Apply the defining matrix action to each input vector. Matrix
    multiplication of a mixed rank-one operator gives the two resulting
    output vectors, with complex conjugation on the bra side.
\end{proof}

\begin{theorem}[Independence of physical output coordinates]
    \label{thm:peps_physical_corrected_basis_invariance}
    \lean{TNLean.PEPS.PairEffect.SourceCircuit.rectangularTraceNorm_physicalCorrectedSourceTerm_basis_eq}
    \leanok
    \uses{thm:peps_prepared_vector_density}
    For a mixed pure input $|\psi\rangle\langle\phi|$, the trace norm of the
    actual physical corrected term is unchanged by replacing either the
    physical or the discard orthonormal basis. The input basis, original
    corrected positions, and arbitrary local correction matrices are fixed.
\end{theorem}

\begin{proof}\leanok
    Expand the actual corrected term in source-basis substitutions.
    Theorem~\ref{thm:peps_prepared_vector_density} expresses every summand as
    the outer product of its actual output vectors. Independence of the
    finite partial trace and trace norm from orthonormal output coordinates
    then applies to their finite sum.
\end{proof}

\begin{definition}[Actual partial Schmidt output]
    \label{def:peps_partial_schmidt_output}
    \lean{TNLean.PEPS.PairEffect.SourceCircuit.partialSchmidtOutput}
    \leanok
    \uses{def:peps_partial_schmidt_assignments}
    For a partial branch $\xi$ and selected endpoint assignment $u$, let
    $\Psi_{\xi,u}(\psi)$ be the vector obtained by preparing the selected
    Schmidt frame columns and the original unselected sources, adjoining the
    owner-relabeled input $\psi$, and applying the canonical partial residual.
    It belongs to the original output memory after the inverse owner
    identification. These are the actual prepared vectors of
    Theorem~\ref{thm:peps_corrected_gaussian_sum}.
\end{definition}

\begin{theorem}[Actual corrected operator on mixed pure inputs]
    \label{thm:peps_gaussian_vector_density}
    \lean{TNLean.PEPS.PairEffect.SourceCircuit.correctedSourceTerm_gaussian_vecMulVec_eq_sum}
    \leanok
    \uses{def:peps_partial_schmidt_output,thm:peps_corrected_gaussian_sum,thm:peps_prepared_vector_density}
    Suppose the affected set meets every corrected source and the weights and
    frames represent every original source vector. For arbitrary $\psi,\phi$,
    \begin{align}
        F_S(\Delta(\omega))(|\psi\rangle\langle\phi|)
        =\sum_{\xi,\zeta}c_{\mathcal A,\xi}\overline{c_{\mathcal A,\zeta}}
        \sum_{u,v}Z^{\xi,\zeta}_{u,v}(\omega)
        |\Psi_{\xi,u}(\psi)\rangle\langle\Psi_{\zeta,v}(\phi)|.
    \end{align}
    The two endpoint assignments are independent, and each partial gate
    coefficient occurs exactly once on each side.
\end{theorem}

\begin{proof}\leanok
    Apply the actual corrected-term matrix expansion and then the
    mixed-vector identity of Theorem~\ref{thm:peps_prepared_vector_density}.
\end{proof}

\begin{theorem}[Square integrability of the actual coefficient products]
    \label{thm:peps_gaussian_products_integrable}
    \lean{TNLean.PEPS.PairEffect.SourceCircuit.sourceCorrection_product_memLp_two}
    \lean{TNLean.PEPS.PairEffect.SourceCircuit.partialSourceCorrection_product_memLp_two}
    \lean{TNLean.PEPS.PairEffect.SourceCircuit.integrable_rectangularTraceNorm_sourceCorrectionSum}
    \leanok
    \uses{thm:peps_actual_gaussian_covariance}
    For $k>0$ and nonnegative local weights, every original-position Gaussian
    coefficient product is square integrable. The same holds after reading
    the local labels from partial branches, provided every selected source
    meets the affected set. Consequently, the trace norm of any finite
    linear combination of these products with fixed rectangular matrix
    coefficients is integrable. Normalization of the weights is unnecessary.
\end{theorem}

\begin{proof}\leanok
    Use the diagonal second moment of the original coefficient
    products. Reindex the retained positions by the original positions.
    Square integrability on the probability space, followed by the finite
    matrix-sum integrability theorem, gives the conclusion.
\end{proof}

\begin{theorem}[Integrability of actual corrected physical operators]
    \label{thm:peps_actual_physical_integrability}
    \lean{TNLean.PEPS.PairEffect.SourceCircuit.integrable_rectangularTraceNorm_map_correctedSourceTerm_gaussian}
    \lean{TNLean.PEPS.PairEffect.SourceCircuit.integrable_rectangularTraceNorm_physicalCorrectedSourceTerm_gaussian}
    \leanok
    \uses{thm:peps_gaussian_products_integrable,thm:peps_corrected_gaussian_sum}
    Let $k>0$, let the weights be nonnegative, and suppose the fixed weights
    and endpoint matrices represent the original source vectors. For any
    fixed complex-linear map $L$ between finite matrix spaces and any input
    matrix $\rho$, the function
    \begin{align}
        \omega\longmapsto\|L(F_S(\Delta(\omega))(\rho))\|_1
    \end{align}
    is integrable whenever the affected set meets $S$. In particular, the
    actual physical corrected term obtained by tracing the final discarded
    registers has integrable trace norm.
\end{theorem}

\begin{proof}\leanok
    The corrected operator is a finite sum of actual coefficient
    products times fixed matrices. Apply
    Theorem~\ref{thm:peps_gaussian_products_integrable}; a fixed linear map
    changes only the matrix coefficients.
\end{proof}

\begin{definition}[A single actual Gaussian branch pair]
    \label{def:peps_gaussian_branch_density}
    \lean{TNLean.PEPS.PairEffect.SourceCircuit.gaussianSchmidtBranchDensity}
    \leanok
    \uses{def:peps_partial_schmidt_output}
    Before multiplying by the outer gate coefficients, define
    \begin{align}
        B_{\xi\zeta}(\omega;\psi,\phi)
        =\sum_{u,v}Z^{\xi,\zeta}_{u,v}(\omega)
        |\Psi_{\xi,u}(\psi)\rangle\langle\Psi_{\zeta,v}(\phi)|.
    \end{align}
    Use finite orthonormal output coordinates for this matrix.
\end{definition}

\begin{theorem}[Original-position expression and branch expectation bound]
    \label{thm:peps_gaussian_branch_density_bound}
    \lean{TNLean.PEPS.PairEffect.SourceCircuit.gaussianSchmidtBranchDensity_eq_sum_original}
    \lean{TNLean.PEPS.PairEffect.SourceCircuit.integrable_rectangularTraceNorm_map_gaussianSchmidtBranchDensity}
    \lean{TNLean.PEPS.PairEffect.SourceCircuit.integral_rectangularTraceNorm_map_correctedSourceTerm_gaussian_le_sum}
    \leanok
    \uses{def:peps_gaussian_branch_density,thm:peps_gaussian_products_integrable,thm:peps_gaussian_vector_density,thm:peps_actual_physical_integrability}
    If the affected set meets $S$, the two sums defining $B_{\xi\zeta}$ may
    be indexed by $\mathcal I_S$ using the actual retained-slot bijection.
    At $e\in S$ the coefficient reads exactly the original local labels
    selected by $\xi,\zeta$; the output vectors are transported by the same
    coordinate bijection.
    For $k>0$ and nonnegative local weights, every
    $\|L(B_{\xi\zeta})\|_1$ is integrable. If the weights and frames also
    represent the original vectors, then
    \begin{align}
        \mathbb E_k\|L(F_S(\Delta)(|\psi\rangle\langle\phi|))\|_1
        \leq\sum_{\xi,\zeta}|c_{\mathcal A,\xi}\overline{c_{\mathcal A,\zeta}}|
        \mathbb E_k\|L(B_{\xi\zeta})\|_1.
    \end{align}
\end{theorem}

\begin{proof}\leanok
    Reindex both finite sums and each source product by the retained
    original positions. Their coefficient products are square integrable.
    Finally use the actual mixed-vector expansion, the trace-norm triangle
    inequality, and linearity of the integral for the resulting finite sum.
\end{proof}

\begin{theorem}[An actual sample bounded by the average error sum]
    \label{thm:peps_sampled_error_integral_bound}
    \lean{TNLean.PEPS.PairEffect.SourceCircuit.exists_sampledSourceMatrix_error_le_sum_integral}
    \leanok
    \uses{thm:peps_actual_physical_integrability,thm:peps_sampled_physical_error}
    Fix $k>0$, nonnegative representing source weights and matrices, finite
    orthonormal input and physical/discard output bases, and an arbitrary
    input matrix. There exists one sample $\omega$ such that
    \begin{align}
        \|\widehat\rho_{\rm phys}(\omega)-\rho_{\rm phys}\|_1
        \leq\sum_{\varnothing\ne S\subseteq\mathcal E}
        \mathbb E_k\|\operatorname{tr}_{\rm disc}F_S(\Delta)\|_1.
    \end{align}
    Here the sampled operator is defined by the actual local source
    replacements, independently of the error expansion.
\end{theorem}

\begin{proof}\leanok
    All summands are nonnegative and integrable. On the probability
    space, their finite sum takes a value no larger than its integral.
    At that sample, apply the established pointwise physical error bound.
\end{proof}

\begin{definition}[Affected parties and retained free sources]
    \label{def:peps_corrected_free_sources}
    \lean{TNLean.PEPS.PairEffect.SourceCircuit.correctedMask}
    \lean{TNLean.PEPS.PairEffect.SourceCircuit.correctedFreeSlots}
    \lean{TNLean.PEPS.PairEffect.SourceCircuit.correctedFreeSlots_eq_true_iff}
    \lean{TNLean.PEPS.PairEffect.SourceCircuit.exists_correctedFreeSlot}
    \leanok
    For $S\subseteq\mathcal E$, let $\mathcal A_S$ be the set of all endpoints
    of sources in $S$. A retained source is free precisely when its original
    position lies in $S$ or its endpoints lie on opposite sides of
    $\mathcal A_S$. Every position in $S$ has its corresponding free retained
    slot. The other retained sources are prepared with their original vectors.
\end{definition}

\begin{theorem}[Actual lifetime and coefficient costs]
    \label{thm:peps_actual_source_lifetime_cost}
    \lean{TNLean.PEPS.PairEffect.SourceCircuit.isTouched_correctedMask_iff}
    \lean{TNLean.PEPS.PairEffect.SourceCircuit.corrected_slot_isTouched}
    \lean{TNLean.PEPS.PairEffect.SourceCircuit.card_touched_correctedMask_le}
    \lean{TNLean.PEPS.PairEffect.SourceCircuit.coefficient_physical_cost_eq}
    \lean{TNLean.PEPS.PairEffect.SourceCircuit.coefficient_physical_cost_le}
    \lean{TNLean.PEPS.PairEffect.SourceCircuit.sum_norm_coefficient_mul_conj_le_lifetime}
    \leanok
    \uses{def:peps_corrected_free_sources}
    A gate occurrence is touched exactly when one of its participating parties
    belongs to $\mathcal A_S$. In particular, the gate containing each source
    in $S$ is touched. Suppose every party in $\mathcal A_S$ participates in
    at most $b$ gate occurrences. Then the number of touched occurrences is
    at most $2b|S|$.
    If each touched gate has coefficient sum at most $B\geq1$, then
    \begin{align}
        \sum_{\xi,\zeta}|c_{\mathcal A_S,\xi}\overline{c_{\mathcal A_S,\zeta}}|
         & \leq B^{4b|S|}.
    \end{align}
    More precisely, this sum times $\prod_{p\in\mathcal A_S}d_p^2$ is the
    corrected-position choice cost for the actual occurrences and their
    branch coefficients. If the prescribed numbers $d_p$ satisfy
    $d_p\leq d$ on $\mathcal A_S$, with $d\geq1$, then
    \begin{align}
        \left(\sum_{\xi,\zeta}|c_{\mathcal A_S,\xi}
        \overline{c_{\mathcal A_S,\zeta}}|\right)
        \prod_{p\in\mathcal A_S}d_p^2
         & \leq (B^{4b}d^4)^{|S|}.
    \end{align}
    The numbers $d_p$ are supplied physical dimensions; no bound on private
    memory is inferred from this assertion.
\end{theorem}

\begin{proof}\leanok
    Every source endpoint participates in its gate. Count touched
    occurrences by summing the lifetime bound over the at most $2|S|$
    affected parties. The exact product formula for partial coefficient
    norms gives the ket and bra factors. Apply the existing corrected-position
    cost estimate, retaining repeated gate occurrences separately.
\end{proof}

\begin{definition}[Two actual contraction factors]
    \label{def:peps_actual_tensor_partition}
    \lean{TNLean.PEPS.PairEffect.Word.HasTensorPartition}
    \leanok
    A word admits a tensor partition for a Boolean owner map when there are
    two allowed source-free words $C,D$, on the respective restricted input
    and output layouts, with $\|C\|,\|D\|\leq1$, such that
    \begin{align}
        P_{\rm out}W=(C\otimes D)P_{\rm in}.
    \end{align}
    The maps $P$ are the canonical owner relabelings followed by the ordered
    register partitions. The equality holds on every input vector.
\end{definition}

\begin{theorem}[Selective preparation gives actual two-side factors]
    \label{thm:peps_corrected_source_partition}
    \lean{TNLean.PEPS.PairEffect.SourceCircuit.exists_corrected_source_partition}
    \leanok
    \uses{def:peps_corrected_free_sources,def:peps_actual_tensor_partition}
    For every allowed chronological composition, every corrected set $S$,
    and every partial branch for $\mathcal A_S$, prepare precisely the
    nonfree sources of Definition~\ref{def:peps_corrected_free_sources} and
    then apply the canonical partial residual. This actual word admits a
    tensor partition into the affected and exterior sides.
    Its free input spaces are fixed across branches. The two factors are
    chosen before any vectors in these free spaces are supplied.
\end{theorem}

\begin{proof}\leanok
    The nonfree source vectors are normalized and both endpoints lie
    on one side. The canonical residual is allowed and source-free.
    The selective-preparation factorization therefore applies on the entire
    free-source memory, with the exterior owners kept together.
\end{proof}

% This common equality entry has the same three owners as the coordinate
% development. On integration, retain its canonical coordinate-chapter home
% and cite that entry here, rather than repeat its declaration tags.
Write $M(a)$ for the tensor memory of a register list $a$.

\begin{theorem}[Head tensor identifications under equality of memories]
    \label{thm:peps_basis_head_memory_equalities}
    \lean{TNLean.PEPS.PairEffect.Layout.selectedHead_heq,
        TNLean.PEPS.PairEffect.Layout.complementaryHead_heq,
        TNLean.PEPS.PairEffect.Layout.add_heq}
    \leanok
    Let $r$ be a register with space $U$, let $a=a'$ and $b=b'$ be
    equal register lists, and suppose $z\in M(a)\otimes M(b)$ and
    $z'\in M(a')\otimes M(b')$ agree after these identifications.
    Write $\alpha$ for the tensor associativity isometry and $L$ for
    the isometry exchanging the first two factors. For every $u\in U$,
    after identifying the memories,
    $\alpha^{-1}(u\otimes z)=\alpha^{-1}(u\otimes z')$ and
    $L(u\otimes z)=L(u\otimes z')$.
    More generally, for equal complex inner product
    spaces $A=B$ and vectors $x,y\in A$ and $x',y'\in B$ agreeing
    under this equality, $x+y$ agrees with $x'+y'$.
\end{theorem}
\begin{proof}\leanok
    Identify the equal lists or spaces and their equal vectors.
    Each pair is then the same tensor or sum.
\end{proof}

\begin{definition}[Concatenating two partitioned layouts]
    \label{def:peps_append_partition}
    \lean{TNLean.PEPS.PairEffect.Layout.appendPartitionIso}
    \lean{TNLean.PEPS.PairEffect.Layout.appendPartitionIso_tmul}
    \leanok
    Given four ordered layouts $a_+,a_-,b_+,b_-$, canonically regroup
    \begin{align}
        (H_{a_+}\otimes H_{a_-})\otimes(H_{b_+}\otimes H_{b_-})
        \simeq H_{a_+b_+}\otimes H_{a_-b_-}.
    \end{align}
    This isometry sends $(x_+\otimes x_-)\otimes(y_+\otimes y_-)$ to
    $(x_+\otimes y_+)\otimes(x_-\otimes y_-)$, with concatenations
    interpreted by their canonical memory isometries.
\end{definition}

\begin{theorem}[Partition commutes with concatenation]
    \label{thm:peps_partition_append}
    \lean{TNLean.PEPS.PairEffect.Layout.partitionIso_append_tmul}
    \leanok
    \uses{def:peps_append_partition}
    Let a Boolean owner map partition two layouts $a,b$. Concatenating their
    memory vectors and then partitioning the resulting layout gives the
    same vector as partitioning each layout first and applying
    Definition~\ref{def:peps_append_partition}. This identity holds for all
    vectors in $H_a$ and $H_b$ and arbitrary owner sets.
\end{theorem}

\begin{proof}\leanok
    \uses{thm:peps_basis_head_memory_equalities}
    Induct on the first layout. For a head register on either side,
    use the defining partition isometry and the four-factor regrouping.
    The empty layout contributes its scalar factor.
\end{proof}

\begin{definition}[An explicit positive sample count]
    \label{def:peps_source_sample_count}
    \lean{TNLean.PEPS.Approximation.sourceSamplingCount}
    \lean{TNLean.PEPS.Approximation.sourceSamplingCount_pos}
    \lean{TNLean.PEPS.Approximation.sourceSamplingCount_lt}
    \leanok
    For $N\geq0$, real $Q$, and $\epsilon>0$, put
    \begin{align}
        k(N,Q,\epsilon)=
        \left\lceil\left(\frac{8\max(1,N)Q}{\epsilon}\right)^2\right\rceil+1.
    \end{align}
    Then $k>0$ and
    $k<(8\max(1,N)Q/\epsilon)^2+2$.
    The ceiling is a nonnegative integer ceiling.
\end{definition}

\begin{theorem}[Finite subset summation and the sample-count estimate]
    \label{thm:peps_source_sample_count_bound}
    \lean{TNLean.PEPS.Approximation.sum_nonemptyFinsets_pow_card}
    \lean{TNLean.PEPS.Approximation.sum_sourceErrorWeights}
    \lean{TNLean.PEPS.Approximation.div_sqrt_sourceSamplingCount_le}
    \lean{TNLean.PEPS.Approximation.sum_sourceErrorWeights_sourceSamplingCount_le}
    \leanok
    \uses{def:peps_source_sample_count}
    For a finite set $\mathcal E$ of cardinality $N$,
    \begin{align}
        \sum_{\varnothing\ne S\subseteq\mathcal E}q^{|S|}
         & =(1+q)^N-1,                            \\
        \sum_{\varnothing\ne S\subseteq\mathcal E}Q^{|S|}k^{-|S|/2}
         & =\left(1+\frac{Q}{\sqrt k}\right)^N-1.
    \end{align}
    The second formula uses $k>0$. If $Q\geq0$ and $\epsilon>0$, the count
    in Definition~\ref{def:peps_source_sample_count} satisfies
    \begin{align}
        \frac{Q}{\sqrt{k(N,Q,\epsilon)}}
         & \leq\frac{\epsilon}{8\max(1,N)}.
    \end{align}
    For $0<\epsilon\leq1$, its finite subset sum is therefore at most
    $\epsilon/4$. These are numerical inequalities; applying them to the
    actual corrected physical terms also requires the corresponding
    individual term estimates.
\end{theorem}

\begin{proof}\leanok
    Expand the product $\prod_{e\in\mathcal E}(1+q)$ and remove the
    empty subset. Substitute $q=Q/\sqrt k$. The ceiling bounds its defining
    square, giving the inverse-square-root estimate. Finally,
    $(1+q)^N\leq\exp(Nq)$ and $Nq\leq\epsilon/8$ imply the stated bound.
    The empty source set gives zero throughout.
\end{proof}

The estimates for specified separated words and the identities for actual
corrected circuit terms have been stated separately. To obtain the full
compression estimate, one must identify the actual free-source input and
physical/discard output registers with the separated words and frames,
then combine the branch, lifetime, and finite-subset estimates.
